\customchapter{INTRODUCCIÓN}

% --- Las tablas de la Introducción se numeran con letras (Tabla A, B, ...)
% porque aquí el contador de capítulo todavía es cero.
\renewcommand{\thetable}{\Alph{table}}

Cuando se levanta una edificación patrimonial con un escáner láser terrestre, el
resultado es una nube de puntos que se escribe como
$P=\{(x_i,f_i)\}_{i=1}^{N}$. Aquí $N$ es la cantidad total de puntos capturados,
que suele ser de millones; $x_i\in\mathbb{R}^3$ son las tres coordenadas del
punto $i$, es decir su ubicación en el espacio; y $f_i$ son los atributos que el
escáner registra junto con esa ubicación, típicamente el color y la intensidad
del retorno del láser. Lo que no aparece en ningún lado es la etiqueta: el
archivo sabe dónde hay materia, pero no sabe qué elemento es cada cosa. Para que
ese levantamiento le sirva a un ingeniero
civil, sea para un análisis estructural, para un modelo BIM o para un informe de
conservación, alguien tiene que asignarle a cada punto la clase del elemento al
que pertenece. Hoy ese trabajo se hace a mano: el especialista recorre el modelo,
selecciona regiones y las va etiquetando una por una. No existe un programa que
lo automatice, y ese vacío es el que este proyecto busca cubrir.

La tarea no se puede resolver bien con fotografías. Un clasificador de imágenes
depende del ángulo de captura y de la iluminación, de modo que la misma columna
vista de frente y vista en escorzo produce descriptores distintos. La nube de
puntos, en cambio, guarda la geometría del elemento sin depender del punto de
vista, que es justamente la propiedad que el problema necesita. Sin embargo, los
métodos de segmentación semántica en tres dimensiones que ya existen tampoco
alcanzan. PTv3~\cite{ptv3} llega a 77,5 mIoU y Sonata~\cite{sonata} a 79,4 mIoU
sobre ScanNet. El mIoU, o \emph{mean intersection over union}, es la métrica
estándar de esta tarea: mide para cada clase qué proporción de la unión entre lo
predicho y lo real quedó bien acertada, y luego promedia ese valor entre todas
las clases; su fórmula se detalla en la Ec.~\eqref{eq:miou}. Ahora bien, los dos
métodos anteriores trabajan con una lista de clases que se fija antes de
entrenar. Cuando aparece un punto que no encaja en ninguna, el modelo no avisa que
no sabe: lo asigna a la clase conocida más parecida. En patrimonio eso es un
problema serio, porque cada edificio tiene elementos propios que no figuran en
ningún catálogo previo.

El enfoque de descubrimiento de clases novedosas, o NCD, elimina ese supuesto y
permite trabajar con clases que el modelo nunca vio etiquetadas. Su mejor método
actual es HyperNCD~\cite{hyperncd}. Este método resume la nube en un conjunto de
\emph{prototipos}, que son vectores representativos alrededor de los cuales se
agrupan los puntos parecidos, y los conecta mediante un \emph{hipergrafo}, una
estructura en la que cada arista puede unir varios prototipos a la vez y no solo
dos, lo que le permite razonar sobre relaciones entre grupos. Cada prototipo se
construye con información geométrica y semántica a la vez.

El problema de HyperNCD es que rinde mal justo donde debería ser fuerte: obtiene entre 11,8 y 47,5 mIoU en las clases
novedosas, frente a 49,8 del modelo supervisado equivalente. Este trabajo
identifica una causa probable de esa debilidad, que es el descriptor geométrico
calculado a mano sobre el que HyperNCD arma sus prototipos, precisamente el tipo
de señal que Sonata señala como perjudicial, y propone reemplazarlo por el
codificador auto-supervisado congelado de Sonata.

El documento se organiza de la siguiente manera. El Capítulo~I presenta la
revisión crítica de la literatura, con el análisis de resultados de las cuatro
fuentes principales, la reproducción de la línea base que fija el punto de
partida y la definición del aporte. El Capítulo~II desarrolla el marco teórico e
incluye las fórmulas de todas las métricas empleadas. El Capítulo~III describe el
marco metodológico, es decir la arquitectura, el punto de fusión, el
procedimiento experimental y los riesgos. El Capítulo~IV reporta los resultados
esperados y su discusión. Al final se presentan las conclusiones y los trabajos
futuros.

\introsection{Formulación del problema}

Dado el levantamiento $P$ de una edificación patrimonial sin etiquetas
semánticas, del que se dispone de un subconjunto etiquetado cuyos puntos
pertenecen a un conjunto de clases conocidas $\mathcal{C}_k$, y de un subconjunto
sin etiquetar cuyos puntos pertenecen a un conjunto de clases novedosas
$\mathcal{C}_n$ que no comparte ninguna clase con el anterior, es decir
$\mathcal{C}_k\cap\mathcal{C}_n=\emptyset$, el problema computacional consiste en
producir una asignación $\hat{y}$ que le entregue a cada punto de $P$ una clase
del conjunto total $\mathcal{C}_k\cup\mathcal{C}_n$, escrito
$\hat{y}:P\rightarrow\mathcal{C}_k\cup\mathcal{C}_n$, de modo que clasifique
correctamente los puntos de las clases conocidas y a la vez descubra
agrupaciones consistentes para los de las clases novedosas. El mejor método
disponible para esta tarea alcanza solo entre 11,8 y 47,5 mIoU sobre $\mathcal{C}_n$ y construye su módulo de prototipos
sobre un descriptor geométrico calculado de forma analítica, con linealidad,
planaridad y dispersión, que es el mismo tipo de señal que la literatura de
aprendizaje auto-supervisado considera un atajo perjudicial. La pregunta de
investigación es entonces si reemplazar ese descriptor por una representación
auto-supervisada confiable aumenta de forma medible el mIoU sobre las clases
novedosas en un levantamiento patrimonial real.

\introsection{Objetivos de investigación}

\paragraph*{Objetivo General}
\addcontentsline{toc}{paragraph}{Objetivo General}

Construir y medir un modelo híbrido que integre las representaciones
auto-supervisadas de Sonata en el módulo de prototipos de HyperNCD, y cuantificar
la mejora del mIoU en la clasificación de elementos estructurales no anotados
sobre el levantamiento de la Catedral de Lima.

\paragraph*{Objetivos Específicos}
\addcontentsline{toc}{paragraph}{Objetivos Específicos}

El proyecto declara \textbf{dos objetivos específicos con trabajo computacional
propio}, OE1 y OE2, cada uno con su resultado verificable y su conclusión
asociada, y \textbf{un objetivo de salida analítica}, OE3.

\medskip
\noindent\textbf{OE1. Integrar Sonata en el módulo de prototipos.} Es el objetivo
central. Cada punto de la nube se describe con un vector $F_i$ que HyperNCD arma
juntando dos partes: $f_{\mathrm{sem}}$, el descriptor semántico que produce su
red supervisada, y $f_{\mathrm{geo}}$, el descriptor geométrico que se calcula a
mano sobre los vecinos del punto (Ec.~\eqref{eq:hyperncd3}). El objetivo consiste
en implementar el bloque de fusión que sustituye ese $f_{\mathrm{geo}}$ por
$f_{\mathrm{ssl}}$, el descriptor que entrega el codificador auto-supervisado de
Sonata con sus pesos congelados, justo antes de la asignación suave a prototipos.
Incluye además resolver la diferencia de resolución entre el codificador PTv3 y
la rejilla de Minkowski mediante \emph{up-cast} y almacenamiento en \emph{cache}
de las características de cada sección.

\begin{itemize}
\item \textit{Herramientas:} pesos públicos de Sonata, Pointcept, PTv3,
\texttt{spconv}, MinkowskiEngine y PyTorch.
\item \textit{Trabajo computacional:} implementación del módulo de fusión,
inferencia congelada de Sonata sobre las cinco secciones y construcción del
hipergrafo de prototipos sobre el descriptor sustituido.
\item \textit{Resultado verificable:} repositorio con el modelo híbrido corriendo
de extremo a extremo y produciendo $F_i=[f_{\mathrm{sem}};f_{\mathrm{ssl}}]$, con
las dimensiones de ambas ramas documentadas y una \textbf{prueba de
funcionamiento} sobre una sección de control, donde la inferencia se completa sin
que el mIoU de las clases conocidas baje respecto de la línea base.
\item \textit{Conclusión asociada:} si la sustitución del descriptor geométrico
por el codificador auto-supervisado es viable dentro del razonamiento por
hipergrafo, y bajo qué condiciones de resolución.
\end{itemize}

\medskip
\noindent\textbf{OE2. Calibrar la configuración del híbrido.} Determinar los
valores de $\alpha$, que balancea la componente geométrica y la semántica, de $M$,
que es el número de vecinos por hiperarista, y de $\epsilon$, que suaviza las
etiquetas, con los que se maximiza el mIoU de las clases novedosas en una sección
de referencia. La calibración se hace teniendo en cuenta que el sistema es
híbrido: al combinar dos métodos, las ventajas y también los modos de falla de
cada uno pasan al modelo resultante, de manera que su comportamiento no se puede
leer como el de ninguno de los dos por separado.

\begin{itemize}
\item \textit{Herramientas:} búsqueda por grilla acotada, AdamW y
Weights~\&~Biases.
\item \textit{Trabajo computacional:} barrido de la grilla sobre la sección de
referencia y fijación de la configuración, que luego se aplica sin reajuste al
resto de las secciones.
\item \textit{Resultado verificable:} curvas de sensibilidad de los tres
parámetros y la configuración fijada, con la métrica reportada también en el
punto intermedio del proceso (Sección~\ref{sec:metricas-hibrido}), lo que permite
atribuir la ganancia a la representación o al razonamiento.
\item \textit{Conclusión asociada:} qué parámetro domina el comportamiento del
híbrido y de qué etapa proviene la mejora.
\end{itemize}

\medskip
\noindent\textbf{OE3. Cuantificar la mejora y delimitar el rango de operación.}
Este objetivo \textbf{no incorpora trabajo computacional nuevo}, porque es la
salida de análisis de OE1 y OE2. Consiste en medir el $\Delta$mIoU de las clases
novedosas (Ec.~\eqref{eq:delta}) del híbrido frente a la línea base reproducida
en la revisión crítica, en descomponer la métrica para identificar qué variable
domina el resultado agregado (Ec.~\eqref{eq:dominante}), en aislar la
contribución del codificador comparando tres variantes del descriptor, que son la
geométrica manual, la auto-supervisada y la concatenación de ambas, y en
determinar de forma empírica la frontera entre el rango favorable y el
desfavorable que se declaran más adelante. Si el análisis muestra que la mejora no
se sostiene, el ajuste se hace sobre OE1 y OE2, no sobre este objetivo.

\begin{itemize}
\item \textit{Herramientas:} el mIoU calculado con emparejamiento húngaro, que es
el procedimiento que decide qué grupo descubierto corresponde a qué clase real
cuando no hay etiquetas; los índices ARI y NMI, que miden si la agrupación
encontrada se parece a la real sin depender de ese emparejamiento; la brecha
cerrada $G$ (Ec.~\eqref{eq:G}), que expresa qué parte de la distancia entre
clases conocidas y novedosas se logró recuperar; y visualizaciones PCA, que
proyectan los descriptores a dos dimensiones para inspeccionarlos a simple vista.
\item \textit{Resultado verificable:} tabla comparativa, estudio de ablación y
curva de degradación que fija el rango declarado.
\end{itemize}

\introsection{Justificación}

La clasificación de elementos estructurales sobre nubes de puntos patrimoniales
se hace hoy a mano porque no hay ninguna herramienta que cubra la tarea cuando la
lista de clases no está cerrada. Cada punto de mIoU que se gana en las clases
novedosas es trabajo de corrección manual que el especialista deja de hacer. Vale
aclarar un matiz importante: los usuarios están conformes con la precisión que
obtienen trabajando a mano, así que el valor de la propuesta no está en igualar
esa precisión sino en automatizar la tarea a un nivel que resulte suficiente. Si
no se demuestra una mejora sobre el estado del arte, la herramienta no tiene
sentido, y por eso el proyecto se plantea como una mejora medible y no como un
desarrollo de software.

Desde el punto de vista computacional, el aporte no es una arquitectura nueva.
Lo que se busca es establecer si un problema que ya fue demostrado en un
escenario con supervisión completa se mantiene en el escenario de descubrimiento
de clases novedosas, donde la calidad de la representación pesa más porque no hay
etiquetas que corrijan al modelo en las clases objetivo. Es un resultado
incierto, medible y que ningún trabajo previo ha reportado, y esa es justamente
la condición que justifica investigarlo. Además, como se mueve una sola pieza del
sistema, cualquier variación en la métrica se puede atribuir al reemplazo y el
estudio de ablación queda sin ambigüedad.

\introsection{Alcance y limitaciones}

\paragraph*{Objeto de estudio}
Se trabaja sobre un único edificio, la \textbf{Catedral de Lima}, con el
levantamiento provisto por la Facultad de Ingeniería Civil. No se evalúan varias
edificaciones para escoger cuál funciona mejor: el objeto es fijo y el modelo se
ajusta a él, dividido en subconjuntos espaciales.

\paragraph*{Generalización declarada}
Los resultados se declaran válidos para estructuras de tipología de catedral o
iglesia colonial, es decir nave con columnas y arcos, cubierta abovedada y
fachada con torres. No se extienden a vivienda, infraestructura vial ni obra
industrial.

\paragraph*{Volumen de datos}
El alcance se declara sobre el levantamiento que está efectivamente disponible y
que se usó en las pruebas preliminares: alrededor de \textbf{1 a 2\,GB} después
de la voxelización. Voxelizar consiste en dividir el espacio en cubos de lado
fijo, llamados vóxeles, y reemplazar todos los puntos que caen dentro de un mismo
cubo por uno solo, de modo que la nube se reduce sin perder la forma del objeto.
Los datos quedan organizados en \textbf{cinco secciones espaciales disjuntas},
que son nave, coro, ábside, torres y fachada, y que funcionan como particiones de
entrenamiento y evaluación. No se declara una escala mayor que no se vaya a
procesar. La Tabla~\ref{tab:datos} detalla la cadena de reducción. El vóxel de
5\,cm es el parámetro que hace tratable el problema y por eso se declara como
límite y no como preferencia. Bajar de ese volumen efectivo le quitaría
significancia al resultado, porque con muy pocos puntos por clase novedosa la
diferencia entre el modelo base y el híbrido dejaría de distinguirse del ruido
entre corridas.

\begin{table}[H]
\centering
\caption{Cadena de reducción de datos sobre el levantamiento real disponible.}
\label{tab:datos}
\begin{tabular}{@{}p{5.6cm}p{3.4cm}p{3.0cm}@{}}
\toprule
\textbf{Etapa} & \textbf{Volumen} & \textbf{Recurso} \\
\midrule
Nube provista (TLS) & 1 a 2\,GB & Almacenamiento \\
Teselado en secciones & 5 particiones disjuntas & CPU \\
Vóxel 5\,cm y normalización & $\sim$2$\times$10$^{7}$ pts & 1 GPU \\
Ventana por paso & $\sim$10$^{5}$ pts & VRAM del nodo \\
Sonata congelado, una pasada con \emph{cache} & Fija & 1 GPU \\
\bottomrule
\end{tabular}
\end{table}

\paragraph*{Hardware}
El proyecto está limitado por cómputo y así se declara. La ejecución se hace
sobre el clúster \textbf{Khipu} de UTEC. La configuración de GPU del clúster
cambió hace poco respecto de la que figuraba en versiones previas de esta
propuesta, de modo que \textbf{no se estiman tiempos de calendario}. En su lugar
se especifica el hardware sobre el que las pruebas se ejecutaron realmente,
verificado por los autores en los nodos que les fueron asignados
(Tabla~\ref{tab:hardware}). El uso de la GPU no se da por viable en una
declaración, sino que se comprueba antes de fijar el tamaño de ventana y el
esquema de \emph{cache}. La división de la catedral en secciones es lo que
sostiene esta política, ya que ninguna etapa necesita la nube completa en
memoria y el experimento sigue siendo ejecutable en una sola GPU dedicada si no
hay asignación multi-GPU disponible.

\begin{table}[H]
\centering
\caption{Requisitos del método y recursos disponibles en el clúster Khipu (UTEC),
verificados por los autores sobre los nodos asignados.}
\label{tab:hardware}
\begin{tabular}{@{}p{2.9cm}p{4.6cm}p{4.5cm}@{}}
\toprule
\textbf{Recurso} & \textbf{Requisito del método} & \textbf{Disponible y verificado} \\
\midrule
GPU & CUDA con capacidad $\geq$ 8.0, una sola unidad; el diseño por secciones evita la necesidad de multi-GPU & 1$\times$ RTX A6000 (\texttt{ds001}, \texttt{g002}); A100 con particiones MIG en \texttt{ag001} \\
VRAM & Por medir en OE3; cota inferior de 0{,}43\,GB solo por los pesos del codificador (108{,}46\,M parámetros) & 48\,GB (49\,140\,MiB) \\
RAM del nodo & Debe alojar la sección completa antes de voxelizar & 64\,GB solicitados por corrida \\
Almacenamiento & \emph{Cache} de $f_{ssl}$, persistente y legible desde ambos entornos & \texttt{/home} compartido, 16\,TB libres sin cuota; \texttt{/local} por nodo \\
\emph{Toolkit} CUDA & 11.3 para MinkowskiEngine (HyperNCD) y 12.4 para PTv3 (Sonata) & 11.4 y 12.4 vía módulos; driver 610.57.04 \\
Conectividad & Descarga de pesos preentrenados y datos & Solo desde el nodo de acceso \\
Gestor de colas & Envío no interactivo de corridas largas & Slurm 25.11.5.1 \\
\bottomrule
\end{tabular}
\end{table}

El nodo \texttt{ag001} ofrece además una NVIDIA A100 particionada en instancias
MIG de 5, 10 y 20\,GB, útiles para pruebas cortas con menor espera en cola. Una
restricción del clúster condiciona el flujo de trabajo y conviene dejarla
registrada: los nodos de cómputo no tienen salida a internet, de modo que los
pesos preentrenados y los datos se descargan desde el nodo de acceso y se leen
después desde el sistema de archivos compartido.

\paragraph*{Costo computacional estimado}
El codificador de Sonata se ejecuta congelado y una sola vez por sección, por lo
que su costo es lineal en el número de puntos retenidos tras la voxelización y no
se repite en cada época de entrenamiento. Para la sección de referencia empleada
en este trabajo, de $7{,}2\times10^{5}$ puntos con normales y color, se estima un
tiempo de inferencia del orden de segundos a decenas de segundos por sección en
la RTX A6000, sin \emph{FlashAttention}. La lectura del archivo ASCII de entrada,
de unos 45\,MB, representa una fracción no despreciable de ese total, por lo que
la conversión previa a formato binario es una optimización disponible si el
cacheo completo de la catedral llegara a ser limitante. Estas cifras son
estimaciones de diseño y serán sustituidas por mediciones directas, que el
procedimiento de cacheo registra por sección.

\paragraph*{Rangos de operación}
El proyecto declara por adelantado dónde se espera que el método funcione y dónde
no, y el objetivo OE3 lo verifica de forma empírica. Se espera que funcione con
densidad efectiva de al menos 1\,pt/cm$^2$ después de la voxelización, con
elementos de 20\,cm o más que tengan una firma geométrica repetida, como
columnas, arcos, bóvedas, cornisas y contrafuertes, con una cobertura de al menos
el 70\% de la superficie del elemento y con hasta unas seis clases novedosas por
sección. Se espera que no funcione bien con ornamento fino de menos de 5\,cm, que
queda por debajo de la resolución del vóxel, con superficies muy ocluidas o
escaneadas desde un solo ángulo, y con clases novedosas que representen menos del
0,5\% de los puntos de la sección, que es el régimen en el que HyperNCD ya
colapsa, con 11,8 mIoU en la partición 3 de SemanticKITTI. Tampoco se espera buen
comportamiento con un número muy alto de clases simultáneas, donde la
representación de Sonata pierde capacidad de distinguirlas, con 29,3 mIoU en
ScanNet200, de 200 clases, frente a 72,5 en ScanNet, de 20 clases, ni con
elementos que se definen más por material o color que por geometría.

\paragraph*{Exclusiones}
Quedan fuera del alcance el preentrenamiento auto-supervisado propio, que en el
caso de Sonata requirió 140 mil escenas y 32 GPU, la detección de instancias, el
vocabulario abierto basado en modelos de lenguaje, la ejecución en tiempo real y
la generación del modelo BIM a partir de la segmentación.

% --- Restauración de la numeración de tablas por capítulo.
\renewcommand{\thetable}{\arabic{chapter}.\arabic{table}}
\setcounter{table}{0}
